\subsection{Self-adjoint extensions of a symmetric operator}

In this section, we propose to classify the self-adjoint extensions of symmetric operators on a complex Hilbert space, and in particular to find conditions ensuring the existence of a self-adjoint extension of such an operator.

Let E be a complex Hilbert space. For every partial operator u on E, we shall write in this section \(u + i\) and \(u - i\) instead of \(u + i1_{E}\) and \(u - i1_{E}\) .

Let u be a closed operator with dense domain on E. We recall (cf. definition 4 of IV, p. 230 and proposition 4 of IV, p. 230) that we denote by \(E_{u}\) the Hilbert space dom(u) equipped with the Hermitian form

\[
(x \mid y) _ {u} = \langle x \mid y \rangle + \langle u (x) \mid u (y) \rangle .
\]

We denote \(\|x\|_{u}\) the norm of an element x of the Hilbert space \(E_{u}\) . Every subspace of dom(u) that is closed in E is closed in \(E_{u}\) (IV, p. 230, remark 1).

DEFINITION 10. --- Let u be a closed symmetric operator on a complex Hilbert space E. Let \(\mathrm{E}_{+} = \mathrm{Ker}(u^{*} - i)\) and \(\mathrm{E}_{-} = \mathrm{Ker}(u^{*} + i)\) . The pair \((\mathrm{E}_{+}, \mathrm{E}_{-})\) of subspaces of \(\mathrm{dom}(u^{*})\) is called the deficiency pair of u. The subspaces \(E_{+}\) and \(E_{-}\) are closed subspaces of E (prop. 3 of IV, p. 229). The pair \((\dim(\mathrm{E}_{+}), \dim(\mathrm{E}_{-}))\) of Hilbertian dimensions of these subspaces is called the deficiency index of u.

PROPOSITION 25. --- Let u be a closed symmetric operator on a complex Hilbert space E. We have

\[
\begin{array}{l} \operatorname{Ker} (u ^ {*} - i) ^ {\circ} = \operatorname{Im} (u + i), \\ \operatorname{Ker} (u ^ {*} + i) ^ {\circ} = \operatorname{Im} (u - i). \end{array}
\]

By prop. 7, c) of IV, p. 236, it suffices to show that the image of \(u + i\) (resp. of \(u - i\) ) is closed in E. Since \(u\) is symmetric, we have \(\langle (u + i)(x) \mid (u + i)(y) \rangle = (x \mid y)_u\) for all \(x\) and \(y\) in \(\text{dom}(u)\) . The map \(x \mapsto u(x) + ix\) from \(E_u\) into E is therefore isometric. Its image is closed in E (lemma 8 of I, p. 107). Since this is also the image of \(u + i\) , the latter is closed in E. Likewise, the image of \(u - i\) is closed in E.

Lemma 7. --- Let u be a closed symmetric operator on a complex Hilbert space E and v a closed symmetric extension of u. The domain of u is a closed subspace of the Hilbert space \(E_{v}\) .

Since v is an extension of u, we have \((x \mid y)_{v} = (x \mid y)_{u}\) for x and y in \(\text{dom}(u)\) . The canonical injection of \(E_{u}\) into \(E_{v}\) is therefore isometric, and the conclusion follows (lemma 8 of I, p. 107).

PROPOSITION 26. --- Let u be a closed symmetric operator on a complex Hilbert space E. Let \((\mathrm{E}_{+}, \mathrm{E}_{-})\) be the deficiency pair of u. The spaces \(E_{+}\) , \(E_{-}\) and dom(u) are closed and mutually orthogonal subspaces of \(E_{u^{*}}\) whose Hilbertian sum is equal to \(E_{u^{*}}\) .

Since u is symmetric, we have \(u \subset u^{*}\) , hence the space \(\text{dom}(u)\) is closed in \(E_{u^{*}}\) (lemma 7). The subspaces \(E_{+}\) and \(E_{-}\) are closed in E and contained in \(\text{dom}(u^{*})\) , hence closed in \(E_{u^{*}}\) .

Let \(x \in \mathrm{E}_{+}\) . We have \(u^{*}(x) = ix\) , whence \(\langle x | u(y) \rangle = -i\langle x | y \rangle\) for all \(y \in \operatorname{dom}(u)\) . Consequently, for every \(y \in \operatorname{dom}(u)\) , we have

\[
(x \mid y) _ {u ^ {*}} = \langle x \mid y \rangle + \langle u ^ {*} (x) \mid u ^ {*} (y) \rangle = i (\langle x \mid u (y) \rangle - \langle x \mid u ^ {*} (y) \rangle) = 0
\]

since u is symmetric. The spaces \(E_{+}\) and \(\operatorname{dom}(u)\) are therefore orthogonal in \(E_{u^{*}}\) . Likewise, \(E_{-}\) and \(\operatorname{dom}(u)\) are orthogonal in \(E_{u^{*}}\) .

Let \(x \in \mathrm{E}_{+}\) and \(y \in \mathrm{E}_{-}\) . Then

\[
(x \mid y) _ {u ^ {*}} = \langle x \mid y \rangle + \langle u ^ {*} (x) \mid u ^ {*} (y) \rangle = \langle x \mid y \rangle - \langle x \mid y \rangle = 0,
\]

hence \(\mathrm{E}_{+}\) and \(\mathrm{E}_{-}\) are orthogonal in \(\mathrm{E}_{u^{*}}\) .

Let \(x \in \mathrm{E}_{u^*}\) orthogonal to the closed subspace \(\operatorname{dom}(u) \oplus \mathrm{E}_+ \oplus \mathrm{E}_-\) . For every \(y \in \operatorname{dom}(u)\) , we have \(u^*(y) = u(y)\) since \(u \subset u^*\) , whence

\[
0 = (x \mid y) _ {u ^ {*}} = \langle x \mid y \rangle + \langle u ^ {*} (x) \mid u ^ {*} (y) \rangle = \langle x \mid y \rangle + \langle u ^ {*} (x) \mid u (y) \rangle .
\]

This implies that \(z = u^{*}(x)\) belongs to the domain of \(u^{*}\) and satisfies \(u^{*}(z) = -x\) . Let \(x_{-} = z - ix\) . It is an element of \(\text{dom}(u^{*})\) satisfying \(u^{*}(x_{-}) = -x - iz = -ix_{-}\) . Hence \(x_{-}\) belongs to \(E_{-}\) . But, for every \(w \in E_{-}\) , we have

\[
\begin{array}{c} - i \langle x _ {-} | w \rangle = - i \langle z | w \rangle + \langle x | w \rangle \\ = \langle u ^ {*} (x) | u ^ {*} (w) \rangle + \langle x | w \rangle = (x | w) _ {u ^ {*}} = 0. \end{array}
\]

Taking \(w = x_{-}\) , we deduce that \(x_{-} = 0\) , that is, \(z = ix\) ; since \(z = u^{*}(x)\) , we therefore have \(x \in \mathrm{E}_{+}\) , whence \(x = 0\) since \(x\) is orthogonal to \(\mathrm{E}_{+}\) in \(\mathrm{E}_{u^{*}}\) . This completes the proof that the direct sum \(\operatorname{dom}(u) \oplus \mathrm{E}_{+} \oplus \mathrm{E}_{-}\) is equal to \(\mathrm{E}_{u^{*}}\) .

COROLLARY. --- Let u be a closed symmetric operator on E. Then u is self-adjoint if and only if the deficiency index of u is equal to (0,0).

Indeed, the symmetric operator \(u\) is self-adjoint if and only if \(\text{dom}(u^*) = \text{dom}(u)\) ; but the proposition shows that \(\text{dom}(u^*)\) is the Hilbert sum of \(\text{dom}(u)\) and of the deficiency subspaces of \(u\) .

\pandocbounded{\includegraphics[keepaspectratio]{images/part002_f39b57bf.jpg}}

Let \(u\) be a closed symmetric partial operator. The spectrum of \(u\) is contained in \(\mathbf{R}\) if and only if \(u\) is self-adjoint. Indeed, we know that \(\mathrm{Sp}(u) \subset \mathbf{R}\) if \(u\) is self-adjoint (prop. 17 of IV, p. 248); if \(u\) is not self-adjoint, one of the subspaces \(\mathrm{Ker}(u^{*} + i)\) or \(\mathrm{Ker}(u^{*} - i)\) is nonzero, hence the image of \(u + i\) or of \(u - i\) is a proper subspace of E (prop. 25 of IV, p. 256), so that \(i\) or \(-i\) belongs to \(\mathrm{Sp}(u)\) .

DEFINITION 11. --- Let u be a closed symmetric operator on a complex Hilbert space E. A boundary condition for u is a partially isometric linear map (EVT, V, p. 41, def. 3) from \(\operatorname{Ker}(u^{*} - i)\) into \(\operatorname{Ker}(u^{*} + i)\) .

Let \(u\) be a closed symmetric operator on a complex Hilbert space E and \(b\) a boundary condition for \(u\) . Let \(I = \text{Ker}(b)^{\circ}\) be the initial subspace of \(b\) . We denote by \(u_b\) the reduction of \(u^*\) to the subspace of \(\text{dom}(u^*)\) that is the direct sum of \(\text{dom}(u)\) and the graph in \(E_+ \oplus E_-\) of the restriction of \(b\) to I. Since \(E_{u^*} = \text{dom}(u) \oplus E_+ \oplus E_-\) (prop. 26), the partial operator \(u_b\) is well defined.

In other words, the domain of \(u_{b}\) is the space of elements \(x \in E\) of the form \(x = x_{0} + y + b(y)\) , where \(x_{0} \in \text{dom}(u)\) and y belongs to the initial subspace of b. We then have \(u_{b}(x_{0} + y + b(y)) = u(x_{0}) + iy - ib(y)\) .

THEOREM 1 (von Neumann). --- Let u be a closed symmetric operator on a complex Hilbert space E and \((\mathrm{E}_{+}, \mathrm{E}_{-})\) its deficiency pair.

\begin{enumerate}
\def\labelenumi{\alph{enumi})}
\item
  For every boundary condition b for u, the partial operator \(u_{b}\) is a closed symmetric extension of u ;
\item
  The map \(b \mapsto u_{b}\) is a bijection from the set of boundary conditions for u onto the set of closed symmetric extensions of u;
\item
  For every boundary condition b for u, we have
\end{enumerate}

\[
\begin{array}{l} \operatorname{Im} (u _ {b} + i) = \operatorname{Im} (u + i) \oplus \operatorname{Im} (b) \\ \operatorname{Im} (u _ {b} - i) = \operatorname{Im} (u - i) \oplus \operatorname{Ker} (b) ^ {\circ}; \end{array}
\]

\begin{enumerate}
\def\labelenumi{\alph{enumi})}
\setcounter{enumi}{3}
\tightlist
\item
  For every boundary condition b for u, the deficiency pair of \(u_{b}\) is \((\operatorname{Ker}(b), \operatorname{Ker}(b^{*}))\) .
\end{enumerate}

Let \(b\) be a boundary condition for \(u\) and \(\mathrm{I} \subset \mathrm{E}_{+}\) its initial subspace. The restriction of \(u^{*}\) to \(\Gamma_{b|I} \subset \mathrm{E}_{+} \oplus \mathrm{E}_{-}\) is the continuous linear map defined by \(x + b(x) \mapsto ix - ib(x)\) for \(x \in I\) . The graph of \(u_{b}\) is the direct sum of the graph of u and the graph of this linear map; it is therefore closed. The closed operator \(u_{b}\) is an extension of u such that \(u_{b} \subset u^{*}\) .

Let \(\gamma_{1}\) and \(\eta_{1}\) be elements of I. Consider the elements

\[
\gamma = \gamma_ {1} + b (\gamma_ {1}), \quad \eta = \eta_ {1} + b (\eta_ {1})
\]

of the graph \(\Gamma_{b|I}\) . We compute

\[
\begin{array}{c} \langle \gamma | u ^ {*} (\eta) \rangle = i \big (\langle b (\gamma_ {1}) | \eta_ {1} \rangle - \langle \gamma_ {1} | b (\eta_ {1}) \rangle \big) + i \big (\langle \gamma_ {1} | \eta_ {1} \rangle - \langle b (\gamma_ {1}) | b (\eta_ {1}) \rangle \big) \\ = i \big (\langle b (\gamma_ {1}) | \eta_ {1} \rangle - \langle \gamma_ {1} | b (\eta_ {1}) \rangle \big) \end{array}
\]

since b is isometric on I. From this we deduce that we have

\[
\langle u ^ {*} (\gamma) \mid \eta \rangle = \langle \gamma \mid u ^ {*} (\eta) \rangle\tag{12}
\]

for all \(\gamma\) and \(\eta\) in \(\Gamma_{b|I}\) .

Let x and y be elements of \(\operatorname{dom}(u)\) , \(\gamma\) and \(\eta\) be elements of \(\Gamma_{b|I}\) . We obtain

\[
\begin{array}{c} \langle x + \gamma | u _ {b} (y + \eta) \rangle = \langle x | u (y) \rangle + \langle x | u ^ {*} (\eta) \rangle + \langle \gamma | u (y) \rangle + \langle \gamma | u ^ {*} (\eta) \rangle \\ = \langle u (x) | y \rangle + \langle u (x) | \eta \rangle + \langle u ^ {*} (\gamma) | y \rangle + \langle u ^ {*} (\gamma) | \eta \rangle \\ = \langle u _ {b} (x + \gamma) | y + \eta \rangle , \end{array}
\]

where we used the fact that u is symmetric and that \(\gamma\) and \(\eta\) belong to the domain of \(u^{*}\) as well as formula (12). The operator \(u_{b}\) is therefore a closed symmetric extension of u. Assertion a) is proved.

The map \(b \mapsto u_b\) is injective. Indeed, a partially isometric map from \(E_+\) into \(E_-\) is uniquely determined by its initial subspace and by its restriction to that subspace. Now the domain of \(u_b\) determines the initial subspace I of \(b\) and the graph of the restriction of \(b\) to I.

Let us prove that the map \(b \mapsto u_b\) is surjective. Let \(w\) be a closed symmetric extension of \(u\) . We have \(w \subset w^* \subset u^*\) . The domain of \(w\) is therefore a subspace of \(\text{dom}(u^*)\) containing \(\text{dom}(u)\) , and \(w\) is the reduction of \(u^*\) to this subspace. Let G be the intersection of the domain of \(w\) and of \(E_+ \oplus E_-\) . This is a closed subspace of \(E_{u^*}\) (Lemma 7) and we have \(\text{dom}(w) = \text{dom}(u) \oplus G\) by proposition 26.

Let \(x \in \mathrm{E}_{+}\) and \(y \in \mathrm{E}_{-}\) such that \(x + y \in \mathrm{G}\) . Since \(w\) is symmetric and is the reduction of \(u^{*}\) to \(\operatorname{dom}(w)\) , we have

\[
\langle x + y \mid u ^ {*} (x + y) \rangle = \langle u ^ {*} (x + y) \mid x + y \rangle .
\]

This equality is equivalent to

\[
\begin{array}{c} \langle x | i x \rangle + \langle x | - i y \rangle + \langle y | i x \rangle + \langle y | - i y \rangle \\ = \langle i x | x \rangle + \langle i x | y \rangle + \langle - i y | x \rangle + \langle - i y | y \rangle , \end{array}
\]

that is, to \(\|x\|^{2} = \|y\|^{2}\) . In particular, the canonical projection \(p_{+}: G \to E_{+}\) is injective, and if I denotes its image, the space G is the graph of an isometric map \(b_{0}\) from I into \(E_{-}\) . For x in I and y in \(E_{-}\) such that \(x + y \in G\) , we have

\[
\begin{array}{r} \| x + y \| _ {u ^ {*}} ^ {2} = \| x \| _ {u ^ {*}} ^ {2} + \| y \| _ {u ^ {*}} ^ {2} = 2 (\| x \| ^ {2} + \| y \| ^ {2}) \\ = 4 \| x \| ^ {2} = 4 \| p _ {+} (x + y) \| ^ {2}. \end{array}
\]

Thus, considering G as a closed subspace of the Hilbert space \(E_{u^{*}}\) , the map \(x \mapsto \frac{1}{2}p_{+}(x)\) is an isometry from G onto I. In particular, I is closed in \(E_{+}\) (lemma 8 of I, p. 107). There then exists a unique boundary condition b for u for which I is the initial space and which coincides with \(b_{0}\) on I. We have \(\text{dom}(w) = \text{dom}(u_{b})\) , whence \(w = u_{b}\) . This establishes assertion b).

Let us prove the assertions \(c\) ) and \(d\) ). Let \(b\) be a boundary condition for \(u\) and I its initial subspace. For \(x_0 \in \mathrm{dom}(u)\) and \(y \in \mathrm{I}\) , we have

\[
u _ {b} (x _ {0} + y + b (y)) - i (x _ {0} + y + b (y)) = (u - i) (x _ {0}) - 2 i b (y).
\]

One has \(\operatorname{Im}(u - i) = \operatorname{E}_{-}^{\circ}\) (prop. 25). Since the image of \(b\) is contained in \(\operatorname{E}_{-}\) , this formula shows that \(\operatorname{Im}(u_b - i) = \operatorname{Im}(u - i) \oplus \operatorname{Im}(b)\) . By loc. cit., we therefore have \(\operatorname{Ker}(u_b^* + i) = \operatorname{E}_{-} \cap \operatorname{Im}(b)^{\circ} = \operatorname{Ker}(b^*)\) .

Similarly, one verifies that \(\operatorname{Im}(u_b + i) = \operatorname{Im}(u + i) \oplus \operatorname{I}\) and hence (loc. cit.) that \(\operatorname{Ker}(u_b^* - i) = \operatorname{I}^\circ \cap \operatorname{E}_+ = \operatorname{Ker}(b)\) .

COROLLARY 1. --- Let u be a closed symmetric operator on a complex Hilbert space E. The map \(b \mapsto u_b\) induces a bijection between the set of isometric isomorphisms from \(\text{Ker}(u^* - i)\) onto \(\text{Ker}(u^* + i)\) and the set of self-adjoint extensions of u. In particular, there exists a self-adjoint extension of u if and only if the two components of the deficiency index of u are equal.

This follows from the corollary of Proposition 26 and the theorem, as well as from EVT, V, p. 25, cor. 2.

COROLLARY 2. --- Let E be a real Hilbert space. Let u be a closed symmetric operator on E. The partial operator \(u_{(\mathbf{C})}\) on \(\mathrm{E}_{(\mathbf{C})}\) is symmetric and closed, and admits at least one self-adjoint extension.

The partial operator \(u_{(\mathbf{C})}\) is closed and symmetric by prop. 2 of IV, p. 228 and lemma 3 of IV, p. 239. The R-linear isomorphism \(j\) from \(\mathrm{E}_{(\mathbf{C})}\) to \(\mathrm{E}_{(\mathbf{C})}\) such that \(j(z \otimes x) = \overline{z} \otimes x\) for all \(z \in \mathbf{C}\) and every \(x \in \mathrm{E}\) induces an isomorphism of Hilbert spaces from \(\mathrm{Ker}(u_{(\mathbf{C})}^* - i)\) onto \(\mathrm{Ker}(u_{(\mathbf{C})}^* + i)\), and we can apply Corollary 1.

COROLLARY 3. --- Let u be a symmetric operator on E. The following conditions are equivalent:

\begin{enumerate}
\def\labelenumi{(\roman{enumi})}
\item
  The operator u is essentially self-adjoint;
\item
  The spaces \(\operatorname{Ker}(u^{*} + i)\) and \(\operatorname{Ker}(u^{*} - i)\) are zero;
\item
  The spaces \(\operatorname{Im}(u + i)\) and \(\operatorname{Im}(u - i)\) are dense in E;
\item
  The partial symmetric operator u has a unique self-adjoint extension.
\end{enumerate}

Conditions (ii) and (iii) are equivalent by Proposition 7, c) of IV, p. 236.

The partial operator \(u\) is essentially self-adjoint if and only if \(\overline{u}\) is self-adjoint, that is, if \(\text{Ker}(\overline{u}^* - i) = \text{Ker}(\overline{u}^* + i) = \{0\}\) (corollary of prop. 26). Since \(\overline{u}^* = u^*\) (prop. 8 of IV, p. 237), condition (i) is therefore equivalent to condition (ii). Moreover, the preceding corollary shows that (ii) implies that \(u\) has a unique self-adjoint extension, which is condition (iv).

Finally, suppose that \(u\) has a unique self-adjoint extension \(v\) . Then the same is true of \(\overline{u}\) , since every self-adjoint extension of \(\overline{u}\) is one of \(u\) , hence is equal to \(v\) . By the preceding corollary, the spaces \(\mathrm{Ker}(\overline{u}^{*} + i)\) and \(\mathrm{Ker}(\overline{u}^{*} - i)\) must be zero, whence condition (ii) using again \(\overline{u}^{*} = u^{*}\) .

Example. --- Let U be an open subset of \(\mathbf{R}^{n}\) endowed with Lebesgue measure \(\mu\) . Let \(\Delta\) be the scalar differential operator

\[
\Delta = - \sum_ {i = 1} ^ {n} \partial_ {i} ^ {2}
\]

on U. Denote by \(u\) the partial operator on the real Hilbert space \(\mathrm{L}_{\mathbf{R}}^{2}(\mathrm{U},\mu)\) with domain \(\mathcal{D}_{\mathbf{R}}(\mathrm{U})\) (IV, p. 201, remark) defined by \(\varphi \mapsto -\sum_{i=1}^{n}\partial_{i}^{2}\varphi\) . We have \(u_{(\mathbf{C})} = \Delta_{-}\) , which is a closable partial operator (prop. 13 of IV, p. 242) and symmetric (since \(\Delta\) is formally symmetric), hence \(u\) is closable and symmetric (proposition 2 of IV, p. 228 and lemma 3 of IV, p. 239). By the above corollary, there therefore exists a self-adjoint extension of \(u_{(\mathbf{C})}\) . This is a Laplacian on U (IV, p. 243).

